% Complete Section 3.4 and Box 3.1, original PDF 37-42 / printed 27-32.
% The floating summary is placed after the final main-text paragraph.
% block: vector-operators-heading page:37 type:section
\subsection{\BT{Example: grad, curl, and div and their adjoints}{Example: grad, curl 和 div 及其伴随}}\label{sec:vector-operators-adjoints}

% block: vector-operators-intro page:37 type:ordinary-paragraph
\BT{In this section we consider the three vector calculus operators which enter the de Rham complex, showing how to view them as closed unbounded operators on Hilbert space, and deriving their basic properties in that framework.}{本节考虑 de Rham 复形中的三个向量微积分算子, 说明如何将它们视为 Hilbert 空间上的闭无界算子, 并在这一框架下推导其基本性质.}

% block: distributional-gradient-domain page:37 type:ordinary-paragraph
\BT{Let $\Omega$ denote a bounded domain in $\mathbb R^n$. We start by considering the unbounded operator from $L^2(\Omega;\mathbb R^n)$ to $L^2(\Omega)$ with domain $C_0^\infty(\Omega;\mathbb R^n)$ given by $\phi\mapsto-\operatorname{div}\phi$ for $\phi\in C_0^\infty(\Omega;\mathbb R^n)$. We shall write this operator as $(-\operatorname{div},C_0^\infty)$ to make clear its domain. Although it is not a closed operator, it is clearly densely defined, so its adjoint $(-\operatorname{div},C_0^\infty)^*$ is a closed operator $L^2(\Omega)\to L^2(\Omega;\mathbb R^n)$. By definition, a real-valued function $u$ on $\Omega$ belongs to the domain of this adjoint if it belongs to $L^2(\Omega)$ and there exists a vector field $v$ in $L^2(\Omega;\mathbb R^n)$ for which}{设 $\Omega$ 表示 $\mathbb R^n$ 中的有界区域. 先考虑从 $L^2(\Omega;\mathbb R^n)$ 到 $L^2(\Omega)$ 的无界算子, 其定义域为 $C_0^\infty(\Omega;\mathbb R^n)$, 对 $\phi\in C_0^\infty(\Omega;\mathbb R^n)$, 由 $\phi\mapsto-\operatorname{div}\phi$ 给出. 为明确其定义域, 将此算子记为 $(-\operatorname{div},C_0^\infty)$. 尽管它不是闭算子, 但显然稠密定义, 故其伴随 $(-\operatorname{div},C_0^\infty)^*$ 是闭算子 $L^2(\Omega)\to L^2(\Omega;\mathbb R^n)$. 按定义, $\Omega$ 上的实值函数 $u$ 属于该伴随的定义域, 是指它属于 $L^2(\Omega)$, 且存在 $L^2(\Omega;\mathbb R^n)$ 中的向量场 $v$, 使得}
\[
-\int_\Omega u\operatorname{div}\phi\,dx=\int_\Omega v\cdot\phi\,dx,\qquad\phi\in C_0^\infty(\Omega;\mathbb R^n).
\]

% block: gradient-hilbert-space page:37 type:ordinary-paragraph
\BT{But this equation is exactly the statement that the distributional gradient of $u$, $\operatorname{grad}u$, belongs to $L^2$ (namely, $\operatorname{grad}u=v=(-\operatorname{div},C_0^\infty)^*u$). Thus the domain of $(-\operatorname{div},C_0^\infty)^*$ is the set of all $L^2$ functions with distributional gradient in $L^2$, which is precisely the definition of the Sobolev space $H^1(\Omega)$. We have shown that}{而此等式恰好表明 $u$ 的分布梯度 $\operatorname{grad}u$ 属于 $L^2$ (即 $\operatorname{grad}u=v=(-\operatorname{div},C_0^\infty)^*u$). 因此, $(-\operatorname{div},C_0^\infty)^*$ 的定义域是所有分布梯度属于 $L^2$ 的 $L^2$ 函数的集合, 这恰好就是 Sobolev 空间 $H^1(\Omega)$ 的定义. 我们已经证明}
\[
(-\operatorname{div},C_0^\infty)^*=(\operatorname{grad},H^1).
\]
\BT{The operator $(\operatorname{grad},H^1)$ not only is closed, but it is evidently also densely defined (its domain contains $C_0^\infty$, for example, which is already dense in $L^2$). The fact that it is closed establishes the completeness of $H^1$.}{算子 $(\operatorname{grad},H^1)$ 不仅是闭的, 而且显然也稠密定义 (例如其定义域包含 $C_0^\infty$, 后者已在 $L^2$ 中稠密). 它是闭算子这一事实确立了 $H^1$ 的完备性.}

% block: divergence-hilbert-space page:37 type:ordinary-paragraph
\BT{We now reverse the roles of div and grad and consider the unbounded operator $(-\operatorname{grad},C_0^\infty)$ from $L^2(\Omega)$ to $L^2(\Omega;\mathbb R^n)$. Reasoning just as above,}{现交换 div 和 grad 的角色, 考虑从 $L^2(\Omega)$ 到 $L^2(\Omega;\mathbb R^n)$ 的无界算子 $(-\operatorname{grad},C_0^\infty)$. 与上面完全相同地推理, 有}
\[
(-\operatorname{grad},C_0^\infty)^*=(\operatorname{div},H(\operatorname{div})),
\]
\BT{where the operator on the right-hand side is closed and densely defined. Its domain is the Hilbert space $H(\operatorname{div};\Omega)$, consisting of $L^2$ vector fields with distributional divergence in $L^2$.}{其中右端算子是闭的且稠密定义. 其定义域为 Hilbert 空间 $H(\operatorname{div};\Omega)$, 由分布散度属于 $L^2$ 的 $L^2$ 向量场组成.}

% block: grad-div-integration-by-parts-intro page:38 type:ordinary-paragraph
\BT{Our next goal is to compute the adjoints of each of the two operators $(\operatorname{grad},H^1)$ and $(\operatorname{div},H(\operatorname{div}))$ just derived. The key to this is the integration-by-parts formula}{下一个目标是计算刚才得到的两个算子 $(\operatorname{grad},H^1)$ 和 $(\operatorname{div},H(\operatorname{div}))$ 各自的伴随. 关键在于分部积分公式}
\begin{equation}\label{eq:grad-div-classical-green}
\int u\operatorname{div}w\,dx+\int\operatorname{grad}u\cdot w\,dx=\int_{\partial\Omega}u|_{\partial\Omega}w|_{\partial\Omega}\cdot n\,ds.
\end{equation}
\BT{This certainly holds for $u\in C^1(\overline\Omega)$, $w\in C^1(\overline\Omega;\mathbb R^n)$ and a smooth boundary, but it is not clear that quantities like the trace of $w$ on $\partial\Omega$ make sense for $w$ only belonging to $H(\operatorname{div})$. So first we investigate that issue. The key result is given in Theorem~\ref{thm:normal-trace-extension}.}{当 $u\in C^1(\overline\Omega)$, $w\in C^1(\overline\Omega;\mathbb R^n)$ 且边界光滑时, 此式当然成立; 但当 $w$ 仅属于 $H(\operatorname{div})$ 时, $w$ 在 $\partial\Omega$ 上的迹这类量是否有意义尚不清楚. 因此先研究这一问题. 关键结论见 Theorem~\ref{thm:normal-trace-extension}.}

% block: sobolev-trace-facts page:38 type:ordinary-paragraph
\BT{For this, and the remainder of the section, we suppose that $\Omega$ has a Lipschitz boundary. We begin by recalling a few elementary facts about the Sobolev space $H^1$. First, the space $C^\infty(\overline\Omega)$ is dense in $H^1(\Omega)$ [\BookNumber{41}, Section 5.3, Theorem 2]. Second, the restriction operator $C^1(\overline\Omega)\to L^2(\partial\Omega)$ extends to a bounded trace operator $\operatorname{tr}:H^1(\Omega)\to L^2(\partial\Omega)$ [\BookNumber{41}, Section 5.5, Theorem 1]. The range of the trace operator is $H^{1/2}(\partial\Omega)$ and, in fact, this is one possible way to define that space. In this approach, the norm in $H^{1/2}(\partial\Omega)$ is given by}{为此, 并在本节余下部分, 假设 $\Omega$ 具有 Lipschitz 边界. 先回顾 Sobolev 空间 $H^1$ 的几个基本事实. 第一, 空间 $C^\infty(\overline\Omega)$ 在 $H^1(\Omega)$ 中稠密 [\BookNumber{41}, Section 5.3, Theorem 2]. 第二, 限制算子 $C^1(\overline\Omega)\to L^2(\partial\Omega)$ 可延拓为有界迹算子 $\operatorname{tr}:H^1(\Omega)\to L^2(\partial\Omega)$ [\BookNumber{41}, Section 5.5, Theorem 1]. 迹算子的像空间为 $H^{1/2}(\partial\Omega)$; 事实上, 这也是定义该空间的一种方式. 按这种方式, $H^{1/2}(\partial\Omega)$ 中的范数由下式给出:}
\begin{equation}\label{eq:trace-quotient-norm}
\|g\|_{1/2,\partial\Omega}=\inf_{\operatorname{tr}u=g}\|u\|_{1,\Omega}.
\end{equation}

% block: trace-annihilator-intro page:38 type:ordinary-paragraph
\BT{Third, the kernel of the trace operator is the space $\mathring H^1(\Omega)$, defined as the closure of $C_0^\infty(\Omega)$ in $H^1(\Omega)$. See [\BookNumber{41}, Section 5.5, Theorem 2]. Knowing both the range and the kernel of the trace operator acting on $H^1(\Omega)$, we can use duality to characterize the functionals on $H^1(\Omega)$ which vanish on the kernel $\mathring H^1(\Omega)$. The following result is elementary.}{第三, 迹算子的核为空间 $\mathring H^1(\Omega)$, 它定义为 $C_0^\infty(\Omega)$ 在 $H^1(\Omega)$ 中的闭包. 见 [\BookNumber{41}, Section 5.5, Theorem 2]. 已知作用于 $H^1(\Omega)$ 的迹算子的像与核后, 便可利用对偶刻画那些在核 $\mathring H^1(\Omega)$ 上为零的 $H^1(\Omega)$ 上的泛函. 下述结论是初等的.}

% block: dual-surjection-annihilator page:38 type:lemma
\Needspace{6\baselineskip}
\begin{Lemma}\label{lem:dual-surjection-annihilator}
\BT{If $\gamma:X\to Y$ is a bounded linear surjection between normed linear spaces with kernel $Z$, the dual map $\gamma':Y'\to X'$ is a bounded injection with range equal to the annihilator of $Z$ (defined as $\{\,L\in X'\mid L|_Z\equiv0\,\}$). Thus, for $L$ in the annihilator there exists a unique $g\in Y'$ such that}{若 $\gamma:X\to Y$ 是赋范线性空间之间的有界线性满射, 核为 $Z$, 则对偶映射 $\gamma':Y'\to X'$ 是有界单射, 其像等于 $Z$ 的零化子 (定义为 $\{\,L\in X'\mid L|_Z\equiv0\,\}$). 因此, 对零化子中的 $L$, 存在唯一 $g\in Y'$, 使得}
\[
Lu=\langle g,\gamma u\rangle,\qquad u\in X.
\]
\end{Lemma}

% block: negative-trace-space-intro page:38 type:ordinary-paragraph
\BT{Applying this lemma with $\gamma$ equal to the trace operator $\operatorname{tr}:H^1(\Omega)\to H^{1/2}(\partial\Omega)$ and defining $H^{-1/2}(\partial\Omega)$ as the dual space of $H^{1/2}(\partial\Omega)$, we obtain the following characterization.}{将此引理用于 $\gamma$ 等于迹算子 $\operatorname{tr}:H^1(\Omega)\to H^{1/2}(\partial\Omega)$ 的情形, 并将 $H^{-1/2}(\partial\Omega)$ 定义为 $H^{1/2}(\partial\Omega)$ 的对偶空间, 得到下述刻画.}

% block: trace-functional-representation page:38 type:proposition
\Needspace{6\baselineskip}
\begin{Proposition}\label{prop:trace-functional-representation}
\BT{If $L:H^1(\Omega)\to\mathbb R$ is a bounded linear functional which annihilates $\mathring H^1(\Omega)$, then there exists a unique $g\in H^{-1/2}(\partial\Omega)$ for which}{若 $L:H^1(\Omega)\to\mathbb R$ 是在 $\mathring H^1(\Omega)$ 上为零的有界线性泛函, 则存在唯一 $g\in H^{-1/2}(\partial\Omega)$, 使得}
\begin{equation}\label{eq:trace-functional-representation}
Lu=\langle g,\operatorname{tr}u\rangle,\qquad u\in H^1(\Omega).
\end{equation}
\BT{Moreover, $\|g\|_{H^{-1/2}(\partial\Omega)}=\|L\|_{H^1(\Omega)'}$.}{此外, $\|g\|_{H^{-1/2}(\partial\Omega)}=\|L\|_{H^1(\Omega)'}$.}
\end{Proposition}

% block: trace-functional-proof page:38 type:proof
\begin{Proof}
\BT{The existence of a unique $g$ satisfying \eqref{eq:trace-functional-representation} comes directly from the lemma. From \eqref{eq:trace-functional-representation} and \eqref{eq:trace-quotient-norm}, we have}{满足 \eqref{eq:trace-functional-representation} 的唯一 $g$ 的存在性直接来自该引理. 由 \eqref{eq:trace-functional-representation} 及 \eqref{eq:trace-quotient-norm}, 有}
\[
|Lu|\leq\|g\|_{H^{-1/2}(\partial\Omega)}\|\operatorname{tr}u\|_{H^{1/2}(\partial\Omega)}\leq\|g\|_{H^{-1/2}(\partial\Omega)}\|u\|_{H^1(\Omega)},
\]
\BT{so $\|L\|\leq\|g\|_{H^{-1/2}(\partial\Omega)}$. For the reverse inequality, let $f\in H^{1/2}(\partial\Omega)$ be arbitrary and choose $u\in H^1(\Omega)$ such that $\operatorname{tr}u=f$ and $\|u\|_1\leq\|f\|_{H^{1/2}(\partial\Omega)}$. Then}{因此 $\|L\|\leq\|g\|_{H^{-1/2}(\partial\Omega)}$. 为证明反向不等式, 任取 $f\in H^{1/2}(\partial\Omega)$, 并选取 $u\in H^1(\Omega)$, 使 $\operatorname{tr}u=f$ 且 $\|u\|_1\leq\|f\|_{H^{1/2}(\partial\Omega)}$. 则}
\[
\langle g,f\rangle=Lu\leq\|L\|\|u\|_1\leq\|L\|\|f\|_{H^{1/2}(\partial\Omega)}.
\]
\end{Proof}

% block: normal-trace-functional page:39 type:ordinary-paragraph
\BT{Next we apply this result to define the trace of $w\in H(\operatorname{div};\Omega)$. The functional}{接下来应用此结论定义 $w\in H(\operatorname{div};\Omega)$ 的迹. 泛函}
\[
L_w(u)=\int u\operatorname{div}w\,dx+\int\operatorname{grad}u\cdot w\,dx,\qquad u\in H^1(\Omega),
\]
\BT{is a bounded linear functional on $H^1$ and, by the definition of the distributional derivative, $L_w(u)$ vanishes for $u\in C_0^\infty(\Omega)$ and hence, by density, for $u\in\mathring H^1(\Omega)$. We may thus apply the proposition to obtain a unique $\gamma_n w\in H^{-1/2}(\partial\Omega)$ for which}{是 $H^1$ 上的有界线性泛函, 且由分布导数的定义, 对 $u\in C_0^\infty(\Omega)$, $L_w(u)$ 为零, 从而由稠密性, 对 $u\in\mathring H^1(\Omega)$ 也为零. 因此可应用该命题得到唯一的 $\gamma_n w\in H^{-1/2}(\partial\Omega)$, 满足}
\begin{equation}\label{eq:normal-trace-green}
\int u\operatorname{div}w\,dx+\int\operatorname{grad}u\cdot w\,dx=\langle\gamma_n w,\operatorname{tr}u\rangle,\quad u\in H^1(\Omega),\ w\in H(\operatorname{div}).
\end{equation}

% block: normal-trace-norm-bound page:39 type:ordinary-paragraph
\BT{Moreover,}{此外,}
\[
\|\gamma_n w\|_{-1/2,\partial\Omega}=\sup_u\frac{\int u\operatorname{div}w\,dx+\int\operatorname{grad}u\cdot w\,dx}{\|u\|_{1,\Omega}}\leq\|w\|_{H(\operatorname{div})}.
\]

% block: normal-trace-classical-identification page:39 type:ordinary-paragraph
\BT{To understand the meaning of $\gamma_n w$, we compare \eqref{eq:normal-trace-green} with the integration-by-parts formula \eqref{eq:grad-div-classical-green}. Fixing $w\in C^1(\overline\Omega;\mathbb R^3)$, note that, as a function of $u\in C^1(\overline\Omega)$, each of the three integrals in \eqref{eq:grad-div-classical-green} is bounded with respect to the $H^1$ norm. The density of smooth functions in $H^1$ then implies \eqref{eq:grad-div-classical-green} for all $u\in H^1(\Omega)$. However, $\gamma_n w$ is uniquely determined by \eqref{eq:normal-trace-green}, so we have}{为理解 $\gamma_n w$ 的含义, 将 \eqref{eq:normal-trace-green} 与分部积分公式 \eqref{eq:grad-div-classical-green} 比较. 固定 $w\in C^1(\overline\Omega;\mathbb R^3)$, 注意到作为 $u\in C^1(\overline\Omega)$ 的函数, \eqref{eq:grad-div-classical-green} 中的三个积分关于 $H^1$ 范数均有界. 由光滑函数在 $H^1$ 中稠密, \eqref{eq:grad-div-classical-green} 便对所有 $u\in H^1(\Omega)$ 成立. 然而 $\gamma_n w$ 由 \eqref{eq:normal-trace-green} 唯一确定, 因此有}
\[
\gamma_n w=w|_{\partial\Omega}\cdot n
\]
\BT{for $w\in C^1(\overline\Omega;\mathbb R^3)$. We have established the following theorem.}{其中 $w\in C^1(\overline\Omega;\mathbb R^3)$. 我们已得到下面的定理.}

% block: normal-trace-extension page:39 type:theorem
\Needspace{6\baselineskip}
\begin{Theorem}\label{thm:normal-trace-extension}
\BT{The operator $w\mapsto w|_{\partial\Omega}\cdot n$ on $C^1(\overline\Omega;\mathbb R^3)$ extends to a bounded linear operator $\gamma_n:H(\operatorname{div})\to H^{-1/2}(\partial\Omega)$. The extension satisfies formula \eqref{eq:normal-trace-green} for integration by parts and the bound $\|\gamma_n w\|_{-1/2,\partial\Omega}\leq\|w\|_{H(\operatorname{div})}$.}{$C^1(\overline\Omega;\mathbb R^3)$ 上的算子 $w\mapsto w|_{\partial\Omega}\cdot n$ 可延拓为有界线性算子 $\gamma_n:H(\operatorname{div})\to H^{-1/2}(\partial\Omega)$. 该延拓满足分部积分公式 \eqref{eq:normal-trace-green} 以及界 $\|\gamma_n w\|_{-1/2,\partial\Omega}\leq\|w\|_{H(\operatorname{div})}$.}
\end{Theorem}

% block: zero-normal-trace-domain page:39 type:ordinary-paragraph
\BT{We are now able to compute the adjoint of the gradient operator with domain $H^1$. For this, define}{现在可以计算定义域为 $H^1$ 的梯度算子的伴随. 为此, 定义}
\[
\mathring H(\operatorname{div})=\{\,w\in H(\operatorname{div})\mid\gamma_n w=0\,\}.
\]

% block: gradient-adjoint page:39 type:theorem
\Needspace{6\baselineskip}
\begin{Theorem}\label{thm:gradient-adjoint-zero-normal}
\BT{The adjoint of the unbounded operator $(\operatorname{grad},H^1)$ is the operator $(-\operatorname{div},\mathring H(\operatorname{div}))$.}{无界算子 $(\operatorname{grad},H^1)$ 的伴随为算子 $(-\operatorname{div},\mathring H(\operatorname{div}))$.}
\end{Theorem}

% block: gradient-adjoint-proof page:39 type:proof
\begin{Proof}
\BT{A vector field $w\in L^2(\Omega;\mathbb R^n)$ is in the domain of the adjoint if and only if there exists $v\in L^2(\Omega)$ such that}{向量场 $w\in L^2(\Omega;\mathbb R^n)$ 属于伴随的定义域, 当且仅当存在 $v\in L^2(\Omega)$, 使得}
\[
\int w\cdot\operatorname{grad}u\,dx=\int vu\,dx,\qquad u\in H^1(\Omega).
\]
\BT{In that case, clearly we must have $w\in H(\operatorname{div})$ and $v=-\operatorname{div}w$ (by considering $u\in C_0^\infty$). Thus $w$ belongs to the domain of the adjoint if and only if $w\in H(\operatorname{div})$ and}{此时, 显然必须有 $w\in H(\operatorname{div})$ 且 $v=-\operatorname{div}w$ (考虑 $u\in C_0^\infty$ 即可). 因而 $w$ 属于伴随的定义域, 当且仅当 $w\in H(\operatorname{div})$ 且}
\[
\int w\cdot\operatorname{grad}u\,dx=-\int(\operatorname{div}w)u\,dx,\qquad u\in H^1(\Omega).
\]
\BT{Comparing with the definition \eqref{eq:normal-trace-green} of $\gamma_n w$, we see that $w$ belongs to the domain of the adjoint if and only if $\gamma_n w=0$.}{与 $\gamma_n w$ 的定义 \eqref{eq:normal-trace-green} 比较, 可见 $w$ 属于伴随的定义域当且仅当 $\gamma_n w=0$.}
\end{Proof}

% block: normal-trace-surjection-intro page:40 type:ordinary-paragraph
\BT{We next show that the normal trace operator maps $H(\operatorname{div})$ not only into $H^{-1/2}(\partial\Omega)$ but onto.}{接下来证明, 法向迹算子不仅将 $H(\operatorname{div})$ 映入 $H^{-1/2}(\partial\Omega)$, 而且映满该空间.}

% block: normal-trace-surjective page:40 type:theorem
\Needspace{6\baselineskip}
\begin{Theorem}\label{thm:normal-trace-surjective}
\BT{The operator $\gamma_n:H(\operatorname{div})\to H^{-1/2}(\partial\Omega)$ is surjective.}{算子 $\gamma_n:H(\operatorname{div})\to H^{-1/2}(\partial\Omega)$ 为满射.}
\end{Theorem}

% block: normal-trace-surjective-proof page:40 type:proof
\begin{Proof}
\BT{Given $g\in H^{-1/2}(\partial\Omega)$, define $u\in H^1$ by the Neumann problem}{给定 $g\in H^{-1/2}(\partial\Omega)$, 通过 Neumann 问题定义 $u\in H^1$:}
\[
\int\operatorname{grad}u\cdot\operatorname{grad}v\,dx+\int uv\,dx=\langle g,\operatorname{tr}v\rangle,\qquad v\in H^1.
\]
\BT{Then $-\Delta u+u=0$. Take $w=\operatorname{grad}u$. Then $\operatorname{div}w=\Delta u=u$, and so $w\in H(\operatorname{div})$. We have}{则 $-\Delta u+u=0$. 取 $w=\operatorname{grad}u$. 则 $\operatorname{div}w=\Delta u=u$, 因此 $w\in H(\operatorname{div})$. 有}
\[
\int w\cdot\operatorname{grad}v\,dx+\int(\operatorname{div}w)v\,dx=\langle g,\operatorname{tr}v\rangle,\qquad v\in H^1,
\]
\BT{which shows that $\gamma_n w=g$.}{这表明 $\gamma_n w=g$.}
\end{Proof}

% block: divergence-adjoint-intro page:40 type:ordinary-paragraph
\BT{Using this, we can argue exactly as in the proof of Theorem~\ref{thm:gradient-adjoint-zero-normal}, to prove the following result, which establishes the adjoint of the divergence operator.}{利用这一点, 可完全按照 Theorem~\ref{thm:gradient-adjoint-zero-normal} 证明中的论证, 证明下述结论, 从而确定散度算子的伴随.}

% block: divergence-adjoint page:40 type:theorem
\Needspace{6\baselineskip}
\begin{Theorem}\label{thm:divergence-adjoint-zero-trace}
\BT{The adjoint of the unbounded operator $(\operatorname{div},H(\operatorname{div}))$ is the operator $(-\operatorname{grad},\mathring H^1)$.}{无界算子 $(\operatorname{div},H(\operatorname{div}))$ 的伴随为算子 $(-\operatorname{grad},\mathring H^1)$.}
\end{Theorem}

% block: reverse-grad-div-adjoints-intro page:40 type:ordinary-paragraph
\BT{Since $T^{**}=T$ for a closed densely defined operator, from Theorems~\ref{thm:gradient-adjoint-zero-normal} and \ref{thm:divergence-adjoint-zero-trace} we can read off two more adjoints.}{由于对闭的稠密定义算子有 $T^{**}=T$, 从 Theorems~\ref{thm:gradient-adjoint-zero-normal} 和 \ref{thm:divergence-adjoint-zero-trace} 可直接读出另外两个伴随.}

% block: reverse-grad-div-adjoints page:40 type:theorem
\Needspace{6\baselineskip}
\begin{Theorem}\label{thm:reverse-gradient-divergence-adjoints}
\[
\begin{gathered}
(\operatorname{div},\mathring H(\operatorname{div}))^*=(-\operatorname{grad},H^1),\\
(\operatorname{grad},\mathring H^1)^*=(-\operatorname{div},H(\operatorname{div})).
\end{gathered}
\]
\end{Theorem}

% block: grad-div-closed-ranges page:40 type:ordinary-paragraph
\BT{Next we verify that all of these operators have closed range. For $(\operatorname{div},H(\operatorname{div}))$ this is clear, since in fact $\operatorname{div}H(\operatorname{div})=L^2$, as can be seen by writing $f\in L^2$ as $\Delta u=\operatorname{div}\operatorname{grad}u$ for some $u\in H^1$. If $\int f\,dx=0$, then we can choose $u$ satisfying the Neumann condition $\operatorname{grad}u\cdot n=0$, and so $f\in\operatorname{div}\mathring H(\operatorname{div})$. Thus the range $\operatorname{div}\mathring H(\operatorname{div})$ has codimension $1$ in $L^2$, and so this is closed as well, by Theorem~\ref{thm:finite-codimension-closed-range}. The closed range theorem then implies that the adjoints of these two maps, namely, $(\operatorname{grad},\mathring H^1)$ and $(\operatorname{grad},H^1)$, both have closed range as well.}{下面验证这些算子的像空间全都是闭的. 对 $(\operatorname{div},H(\operatorname{div}))$, 这一点很明显, 因为实际上 $\operatorname{div}H(\operatorname{div})=L^2$: 将 $f\in L^2$ 写成某个 $u\in H^1$ 的 $\Delta u=\operatorname{div}\operatorname{grad}u$ 即可看出. 若 $\int f\,dx=0$, 则可选取 $u$ 满足 Neumann 条件 $\operatorname{grad}u\cdot n=0$, 从而 $f\in\operatorname{div}\mathring H(\operatorname{div})$. 因此像空间 $\operatorname{div}\mathring H(\operatorname{div})$ 在 $L^2$ 中的余维为 $1$, 所以由 Theorem~\ref{thm:finite-codimension-closed-range}, 它也是闭的. 闭值域定理进而表明这两个映射的伴随, 即 $(\operatorname{grad},\mathring H^1)$ 和 $(\operatorname{grad},H^1)$, 也都具有闭像空间.}

% block: curl-integration-by-parts-intro page:40 type:ordinary-paragraph
\BT{Finally, we turn to the curl operator, assuming the domain $\Omega\subset\mathbb R^3$. Most of the arguments we applied above to grad and div apply to curl as well, but not all, because neither the kernel nor the range of curl is finite codimensional. For the curl, the relevant integration-by-parts formula is}{最后转向 curl 算子, 假设区域 $\Omega\subset\mathbb R^3$. 上面对 grad 和 div 使用的大多数论证也适用于 curl, 但并非全部适用, 因为 curl 的核与像都不是有限余维的. 对 curl, 相应的分部积分公式为}
\begin{equation}\label{eq:curl-classical-green}
\begin{aligned}
\int v\cdot\operatorname{curl}u\,dx-\int\operatorname{curl}v\cdot u\,dx&=\int_{\partial\Omega}v|_{\partial\Omega}\times n\cdot u|_{\partial\Omega}\,ds,\\
&\hspace{-7em}u,v\in C^1(\overline\Omega;\mathbb R^3),
\end{aligned}
\end{equation}

% block: curl-distributional-definition pages:40,41 type:ordinary-paragraph
\BT{where $n$ is the unit outward normal (which is defined almost everywhere since the boundary is Lipschitz). The space $H(\operatorname{curl})$ is defined to be the space of vector fields in $L^2(\Omega;\mathbb R^3)$ whose curl taken in the sense of distributions also belongs to $L^2(\Omega;\mathbb R^3)$. Thus, a vector field $u\in L^2(\Omega;\mathbb R^3)$ belongs to $H(\operatorname{curl})$ if and only if there exists a vector field $v\in L^2(\Omega;\mathbb R^3)$ such that}{其中 $n$ 为单位外法向量 (由于边界为 Lipschitz, 它几乎处处有定义). 空间 $H(\operatorname{curl})$ 定义为那些分布意义下的旋度也属于 $L^2(\Omega;\mathbb R^3)$ 的 $L^2(\Omega;\mathbb R^3)$ 向量场所成的空间. 因此, 向量场 $u\in L^2(\Omega;\mathbb R^3)$ 属于 $H(\operatorname{curl})$, 当且仅当存在向量场 $v\in L^2(\Omega;\mathbb R^3)$, 使得}
\begin{equation}\label{eq:distributional-curl}
\int_\Omega u\cdot\operatorname{curl}w\,dx=\int_\Omega v\cdot w\,dx,\qquad w\in C_0^\infty(\Omega;\mathbb R^3),
\end{equation}
\BT{and, in this case, $v=\operatorname{curl}u$. Thus, just as above,}{且此时 $v=\operatorname{curl}u$. 因此, 与上面相同,}
\[
(\operatorname{curl},C_0^\infty)^*=(\operatorname{curl},H(\operatorname{curl})).
\]
\BT{Being an adjoint, this last is a closed operator, which means that $H(\operatorname{curl})$ is a Hilbert space when endowed with the graph norm}{作为伴随, 后一个算子是闭算子; 这意味着 $H(\operatorname{curl})$ 配备如下图范数后是 Hilbert 空间:}
\[
\|u\|_{H(\operatorname{curl})}^2=\|u\|_{L^2}^2+\|\operatorname{curl}u\|_{L^2}^2.
\]

% block: tangential-trace-functional page:41 type:ordinary-paragraph
\BT{Now we develop a trace for $v\in H(\operatorname{curl})$. The map}{现在为 $v\in H(\operatorname{curl})$ 建立迹. 映射}
\[
u\mapsto\int v\cdot\operatorname{curl}u\,dx-\int\operatorname{curl}v\cdot u\,dx,\qquad u\in H^1(\Omega;\mathbb R^3),
\]
\BT{is a bounded linear functional on $H^1(\Omega;\mathbb R^3)$ which vanishes on $C_0^\infty(\Omega;\mathbb R^3)$ and so, by density, on $\mathring H^1(\Omega;\mathbb R^3)$. Therefore we can apply Proposition~\ref{prop:trace-functional-representation} to infer the existence of a unique $\gamma_\tau v\in H^{-1/2}(\partial\Omega;\mathbb R^3)$ for which}{是 $H^1(\Omega;\mathbb R^3)$ 上的有界线性泛函, 在 $C_0^\infty(\Omega;\mathbb R^3)$ 上为零, 因而由稠密性, 在 $\mathring H^1(\Omega;\mathbb R^3)$ 上也为零. 因此可应用 Proposition~\ref{prop:trace-functional-representation}, 推出存在唯一的 $\gamma_\tau v\in H^{-1/2}(\partial\Omega;\mathbb R^3)$, 使得}
\begin{equation}\label{eq:tangential-trace-green}
\begin{aligned}
\int v\cdot\operatorname{curl}u\,dx-\int\operatorname{curl}v\cdot u\,dx&=\langle\gamma_\tau v,\operatorname{tr}u\rangle,\\
u\in H^1(\Omega;\mathbb R^3),\quad v&\in H(\operatorname{curl}),
\end{aligned}
\end{equation}
\BT{and to obtain the norm bound}{并得到范数界}
\begin{equation}\label{eq:tangential-trace-bound}
\|\gamma_\tau v\|_{-1/2,\partial\Omega}\leq\|v\|_{H(\operatorname{curl})}.
\end{equation}

% block: tangential-trace-identification page:41 type:ordinary-paragraph
\BT{Comparing with \eqref{eq:curl-classical-green} we see that we have extended the \emph{tangential trace operator} $\gamma_\tau v:=v|_{\partial\Omega}\times n$ to a bounded linear map on $H(\operatorname{curl})$ with range belonging to the Sobolev space of $H^{-1/2}(\partial\Omega;\mathbb R^3)$ of vector fields on $\partial\Omega$, and so we have established the analogue of Theorem~\ref{thm:normal-trace-extension} for $H(\operatorname{curl})$.}{与 \eqref{eq:curl-classical-green} 比较可见, 我们已将\emph{切向迹算子} $\gamma_\tau v:=v|_{\partial\Omega}\times n$ 延拓为 $H(\operatorname{curl})$ 上的有界线性映射, 其像属于 $\partial\Omega$ 上向量场的 Sobolev 空间 $H^{-1/2}(\partial\Omega;\mathbb R^3)$. 因而得到了 $H(\operatorname{curl})$ 的相应于 Theorem~\ref{thm:normal-trace-extension} 的结论.}

% block: tangential-trace-extension page:41 type:theorem
\Needspace{6\baselineskip}
\begin{Theorem}\label{thm:tangential-trace-extension}
\BT{The operator $v\mapsto v|_{\partial\Omega}\times n$ on $C^1(\overline\Omega;\mathbb R^3)$ extends to a bounded linear operator $\gamma_\tau:H(\operatorname{curl})\to H^{-1/2}(\partial\Omega;\mathbb R^3)$. The extension satisfies the integration-by-parts formula \eqref{eq:tangential-trace-green} and the bound \eqref{eq:tangential-trace-bound}.}{$C^1(\overline\Omega;\mathbb R^3)$ 上的算子 $v\mapsto v|_{\partial\Omega}\times n$ 可延拓为有界线性算子 $\gamma_\tau:H(\operatorname{curl})\to H^{-1/2}(\partial\Omega;\mathbb R^3)$. 此延拓满足分部积分公式 \eqref{eq:tangential-trace-green} 和界 \eqref{eq:tangential-trace-bound}.}
\end{Theorem}

% block: tangential-trace-not-surjective page:41 type:ordinary-paragraph
\BT{Note, however, that the trace operator $\gamma_\tau$ is not onto $H^{-1/2}(\partial\Omega;\mathbb R^3)$, i.e., the analogue of Theorem~\ref{thm:normal-trace-surjective} does not hold. For one thing, $\gamma_\tau w$ is always a tangential vector field. This is clear at least if $v$ and $\partial\Omega$ are smooth, while the precise way to define this in the case of a merely Lipschitz boundary is subtle. There are further restrictions on the range of $\gamma_\tau$ as well. See [\BookNumber{25}, \BookNumber{26}] for details.}{然而须注意, 迹算子 $\gamma_\tau$ 并不映满 $H^{-1/2}(\partial\Omega;\mathbb R^3)$, 即相应于 Theorem~\ref{thm:normal-trace-surjective} 的结论不成立. 首先, $\gamma_\tau w$ 总是切向向量场. 至少在 $v$ 和 $\partial\Omega$ 光滑时, 这一点很清楚, 而在边界仅为 Lipschitz 的情形, 如何精确定义这一点却很微妙. $\gamma_\tau$ 的像还受到其他限制. 详情参见 [\BookNumber{25}, \BookNumber{26}].}

% block: curl-adjoint-calculation pages:41,42 type:ordinary-paragraph
\BT{Finally, we compute the adjoint of $(\operatorname{curl},H(\operatorname{curl}))$. If $w$ belongs to the domain of the adjoint, then there exists $v=(\operatorname{curl},H(\operatorname{curl}))^*w\in L^2(\Omega;\mathbb R^3)$ with}{最后计算 $(\operatorname{curl},H(\operatorname{curl}))$ 的伴随. 若 $w$ 属于伴随的定义域, 则存在 $v=(\operatorname{curl},H(\operatorname{curl}))^*w\in L^2(\Omega;\mathbb R^3)$, 满足}
\[
\int w\cdot\operatorname{curl}u\,dx=\int v\cdot u\,dx,\qquad u\in H(\operatorname{curl}).
\]
\BT{Such a $w$ must belong to $H(\operatorname{curl})$ and have $\operatorname{curl}w=v$. Thus $w$ belongs to the domain of $(\operatorname{curl},H(\operatorname{curl}))^*$ if and only if $w\in H(\operatorname{curl})$ and}{这样的 $w$ 必须属于 $H(\operatorname{curl})$, 且 $\operatorname{curl}w=v$. 因此 $w$ 属于 $(\operatorname{curl},H(\operatorname{curl}))^*$ 的定义域, 当且仅当 $w\in H(\operatorname{curl})$ 且}
\[
\int w\cdot\operatorname{curl}u\,dx=\int\operatorname{curl}w\cdot u\,dx,\qquad u\in H(\operatorname{curl}).
\]
\BT{It can be shown (by a mollification argument similar to that for $H^1$) that smooth functions are dense in $H(\operatorname{curl})$. Thus we need only require that the last equation hold for all $u\in H^1(\Omega;\mathbb R^3)$. But then \eqref{eq:tangential-trace-green} implies that this holds if and only if $\gamma_\tau w=0$, and so we have established that the adjoint of $(\operatorname{curl},H(\operatorname{curl}))$ is $(\operatorname{curl},\mathring H(\operatorname{curl}))$. Finally, we can take adjoints and get the reverse relation. These results are summarized in the following theorem.}{可以证明 (用与 $H^1$ 情形类似的磨光论证), 光滑函数在 $H(\operatorname{curl})$ 中稠密. 因此只需要求上一个等式对所有 $u\in H^1(\Omega;\mathbb R^3)$ 成立. 而由 \eqref{eq:tangential-trace-green}, 这当且仅当 $\gamma_\tau w=0$ 时成立, 因而我们证明了 $(\operatorname{curl},H(\operatorname{curl}))$ 的伴随是 $(\operatorname{curl},\mathring H(\operatorname{curl}))$. 最后取伴随, 可得到反向关系. 这些结论概括在下面的定理中.}

% block: curl-adjoints page:42 type:theorem
\Needspace{6\baselineskip}
\begin{Theorem}\label{thm:curl-adjoints}
\[
\begin{gathered}
(\operatorname{curl},H(\operatorname{curl}))^*=(\operatorname{curl},\mathring H(\operatorname{curl})),\\
(\operatorname{curl},\mathring H(\operatorname{curl}))^*=(\operatorname{curl},H(\operatorname{curl})),
\end{gathered}
\]
\BT{where}{其中}
\[
\mathring H(\operatorname{curl}):=\{\,w\in H(\operatorname{curl})\mid\gamma_\tau w=0\,\}.
\]
\end{Theorem}

% block: curl-closed-range-forward-reference page:42 type:ordinary-paragraph
\BT{In Section~\ref{sec:l2-de-rham-complex} we show that the range $\operatorname{curl}H(\operatorname{curl})$ is closed.}{在 Section~\ref{sec:l2-de-rham-complex} 中, 我们将证明像空间 $\operatorname{curl}H(\operatorname{curl})$ 是闭的.}

% block: unbounded-operators-summary-heading page:42 type:ordinary-paragraph
\begin{tcolorbox}[enhanced,breakable,colback=KuangNei,frame hidden,boxrule=0pt,sharpish corners,before skip=10pt,after skip=10pt]
\refstepcounter{booksummary}\label{box:unbounded-operators-summary}
\noindent\textbf{Box~\thebooksummary:} \BT{Summary of Chapter~\ref{chap:unbounded-operators} on operators on Hilbert spaces.}{Chapter~\ref{chap:unbounded-operators} 中 Hilbert 空间上算子的内容摘要.}\par

% block: summary-unbounded-domain page:42 type:ordinary-paragraph
\BT{An \emph{(unbounded) operator} from one Hilbert space $X$ to another $Y$ is a linear operator mapping its domain $D(T)$, a subspace of $X$, into $Y$. The domain need not be closed in $X$ and the operator need not be bounded, but usually we require that the operator be \emph{densely defined}, i.e., that $D(T)$ is dense in $X$.}{从 Hilbert 空间 $X$ 到另一个空间 $Y$ 的\emph{(无界) 算子}, 是将其定义域 $D(T)$ (它是 $X$ 的子空间) 映入 $Y$ 的线性算子. 定义域不必在 $X$ 中闭, 算子也不必有界, 但通常要求算子\emph{稠密定义}, 即 $D(T)$ 在 $X$ 中稠密.}

% block: summary-closed-operator page:42 type:ordinary-paragraph
\BT{An operator is \emph{closed} if its graph is closed in $X\times Y$ or, equivalently, if $D(T)$ is complete when furnished with the \emph{graph norm}. An everywhere defined operator is closed if and only if it is bounded.}{若算子的图在 $X\times Y$ 中闭, 或等价地, $D(T)$ 配备\emph{图范数}后完备, 则称该算子是\emph{闭的}. 处处有定义的算子是闭的当且仅当它有界.}

% block: summary-adjoint-operator page:42 type:ordinary-paragraph
\BT{A densely defined operator $T$ from $X$ to $Y$ gives rise to an adjoint operator $T^*$ from $Y$ to $X$. The adjoint $T^*$ is always closed. If $T$ is closed, then $T^*$ is also densely defined, and, in this case, $T^{**}=T$.}{从 $X$ 到 $Y$ 的稠密定义算子 $T$ 产生从 $Y$ 到 $X$ 的伴随算子 $T^*$. 伴随 $T^*$ 总是闭的. 若 $T$ 是闭的, 则 $T^*$ 也稠密定义, 并且此时 $T^{**}=T$.}

% block: summary-closed-range page:42 type:ordinary-paragraph
\BT{The range of a closed, densely defined linear operator $X\to Y$ need not be a closed subspace, but the range of $T$ is closed in $Y$ if and only if the range of $T^*$ is closed in $X$. If $\dim Y/\mathcal R(T)<\infty$, then $T$ has closed range. If the inclusion $D(T)\subset X$ is compact, then $T$ has closed range.}{闭的稠密定义线性算子 $X\to Y$ 的像空间不一定是闭子空间, 但 $T$ 的像空间在 $Y$ 中闭当且仅当 $T^*$ 的像空间在 $X$ 中闭. 若 $\dim Y/\mathcal R(T)<\infty$, 则 $T$ 具有闭像空间. 若包含映射 $D(T)\subset X$ 是紧的, 则 $T$ 具有闭像空间.}

% block: summary-vector-operators page:42 type:ordinary-paragraph
\BT{On a bounded domain in $\mathbb R^3$ with Lipschitz boundary, the operators grad, curl, and div may be viewed as unbounded operators on $L^2$ (scalars or vectors), with domains $H^1$, $H(\operatorname{curl})$, $H(\operatorname{div})$, respectively. They are densely defined, closed, and have closed ranges. Their adjoints are $-\operatorname{div}$, $\operatorname{curl}$, and $-\operatorname{grad}$ with domains $\mathring H(\operatorname{div})$, $\mathring H(\operatorname{curl})$, and $\mathring H^1$, respectively.}{在 $\mathbb R^3$ 中具有 Lipschitz 边界的有界区域上, 算子 grad, curl 和 div 可视为 $L^2$ (标量或向量) 上的无界算子, 定义域分别为 $H^1$, $H(\operatorname{curl})$ 和 $H(\operatorname{div})$. 它们稠密定义, 闭, 且具有闭像空间. 它们的伴随分别为 $-\operatorname{div}$, $\operatorname{curl}$ 和 $-\operatorname{grad}$, 定义域分别为 $\mathring H(\operatorname{div})$, $\mathring H(\operatorname{curl})$ 和 $\mathring H^1$.}
\end{tcolorbox}
